% 第 5 章：函数的概念与基本初等函数（集合结构视角）
\chapter{函数的概念与基本初等函数（集合结构视角）}

\section{基础·三视角互译}

\begin{problem}\level{易}
已知 $A=\{1,2\}$，$B=\{2,3,4\}$，设 $R=\{(1,2),(1,3),(2,4)\}\subseteq A\times B$。
(1) 判断 $R$ 是否为从 $A$ 到 $B$ 的函数，说明理由；
(2) 在 $A\times B$ 中另写出一个“一对多”的反例关系；
(3) 求从 $A$ 到 $B$ 的函数总个数。
\end{problem}
\begin{solution}
\paragraph{集合与点集本质剖析}
函数 $f:A\to B$ 是 $A\times B$ 的满足“每个 $x\in A$ 恰出现一次作为第一坐标”的子集，即点集 $\{(x,f(x))\mid x\in A\}$ 中任意两点的横坐标都不相同。$R$ 中 $(1,2)$ 与 $(1,3)$ 的横坐标都是 $1$，唯一性被破坏，所以 $R$ 不是函数。
\paragraph{标准解题步骤}
先列出 $A\times B=\{(1,2),(1,3),(1,4),(2,2),(2,3),(2,4)\}$，共 $6$ 个有序数对。检验 $R$：$x=1$ 同时对应 $2$ 和 $3$ 两个象，违反单值性，故 $R$ 不是函数。反例可取 $R'=\{(1,2),(1,3),(2,2)\}$，其中 $x=1$ 对应两个值。计数时，每个 $x\in A$ 都要在 $B$ 中选定唯一象，各有 $3$ 种选择，故函数总数为 $3^2=9$；若要求一一映射，则第一个 $x$ 有 $3$ 种、第二个 $x$ 有 $2$ 种，共 $3\times2=6$ 个。
\paragraph{学生典型易错点诊断}
最易犯的错误是把“存在 $y$”误当成“存在唯一 $y$”，看到 $R$ 覆盖了 $A$ 中每个元素就判它是函数。二是把 $A\times B$ 本身当成函数，实际上它含有大量一对多的数对。三是计数时忽略“每个 $x$ 都必须有象”，错算成 $2^3=8$；四是把函数个数与单射个数混为一谈，$9$ 是全体函数数，$6$ 才是一一映射数。规避方法：判函数先查“同一 $x$ 是否只出现一次”，计数时用分步乘法原理逐 $x$ 处理。
\end{solution}

\begin{problem}\level{易}
设 $A=\{-1,0,1\}$，$B=\{1,2,3\}$，映射 $f:A\to B$ 定义为 $f(x)=x+2$。
(1) 用列举法写出 $f$ 的有序数对集合；
(2) 在平面直角坐标系中描出这些点，并说明能否把它们连成一条连续曲线；
(3) 画出从 $A$ 到 $B$ 的箭头示意图；
(4) 判断 $f$ 是否为一一映射，若是，写出逆映射 $f^{-1}$ 的对应关系。
\end{problem}
\begin{solution}
\paragraph{集合与点集本质剖析}
同一个映射有三种等价语言：数对语言 $G_f=\{(x,f(x))\}\subseteq A\times B$，几何语言即坐标系中的点集，对应语言即从 $A$ 到 $B$ 的箭头图。三者信息完全一致，只是呈现方式不同；本题的核心就是把三种语言互译。
\paragraph{标准解题步骤}
$f(-1)=1$，$f(0)=2$，$f(1)=3$，故 $G_f=\{(-1,1),(0,2),(1,3)\}$。在坐标系中描出 $(-1,1),(0,2),(1,3)$，它们都在直线 $y=x+2$ 上；但 $A$ 是离散点集而非区间，只能描点，不能连线成连续曲线。箭头示意图如下：
\begin{center}
\begin{tikzpicture}[>=stealth]
\node (a) at (0,1.5) {$-1$}; \node (b) at (0,0) {$0$}; \node (c) at (0,-1.5) {$1$};
\node (d) at (3,1.5) {$1$}; \node (e) at (3,0) {$2$}; \node (f) at (3,-1.5) {$3$};
\draw[->] (a)--(d); \draw[->] (b)--(e); \draw[->] (c)--(f);
\node[above] at (0,2) {$A$}; \node[above] at (3,2) {$B$};
\end{tikzpicture}
\end{center}
由于 $B$ 中每个元素都恰有一个原像，$f$ 是一一映射；逆映射为 $f^{-1}(1)=-1$，$f^{-1}(2)=0$，$f^{-1}(3)=1$，即 $f^{-1}(x)=x-2\ (x\in B)$。
\paragraph{学生典型易错点诊断}
最易犯的错误是把定义域 $A=\{-1,0,1\}$ 误当成区间 $[-1,1]$，从而把三个点用线连起来，混淆了“离散点集”与“连续曲线”。二是在箭头图中画出一对多的箭头，或让 $B$ 中元素没有原像，误以为映射不必覆盖 $B$。三是写逆映射时忘记交换定义域与值域，把 $f^{-1}$ 的定义域写成 $\mathbb{R}$。规避方法：始终先问“定义域到底是哪个集合”，再在三种语言间逐条核对。
\end{solution}

\begin{problem}\level{中}
用垂线检验法判断下列点集或曲线中哪些是函数图象，哪些不是，并指出破坏函数性的点：
(1) $S=\{(1,2),(1,-2),(2,3),(3,0)\}$；
(2) 抛物线 $y=x^2$；
(3) 曲线 $y^2=x\ (x\geq0)$；
(4) 竖直线 $x=1$。
\end{problem}
\begin{solution}
\paragraph{集合与点集本质剖析}
垂线检验法是单值性的几何化身：任意竖直线 $x=a$ 与图象至多有一个交点，等价于每个 $x$ 至多对应一个 $y$。它把“函数定义”翻译成“点集的交点个数”问题，是本层最核心的观点。
\paragraph{标准解题步骤}
(1) 直线 $x=1$ 与 $S$ 交于 $(1,2)$ 和 $(1,-2)$ 两点，故 $S$ 不是函数图象，破坏点是 $x=1$。(2) 任意竖直线与抛物线 $y=x^2$ 恰有一个交点，故是函数图象。(3) 直线 $x=1$ 与曲线交于 $(1,1)$ 和 $(1,-1)$，一对多，故不是函数；它实际是两支 $y=\pm\sqrt{x}$ 的并。(4) 竖直线 $x=1$ 与自身重合，交点无穷多，当然不是函数。
\paragraph{学生典型易错点诊断}
最易犯的错误是凭“看起来像曲线”就判为函数，忽视竖直线检验；二是处理 $y^2=x$ 时只写 $y=\sqrt{x}$ 而丢掉 $y=-\sqrt{x}$，误判为函数；三是把“一个 $y$ 对应多个 $x$”（如 $y=x^2$ 中 $y=4$ 对应 $\pm2$，这是允许的）与“一个 $x$ 对应多个 $y$”（禁止）混为一谈。规避方法：判函数时只数“同一竖直线上的点”，看到竖直线直接判否。
\end{solution}

\section{进阶·点集交集与区域}

\begin{problem}\level{中}
设 $f(x)=x^2-2x$，$g(x)=x$。
(1) 求方程 $f(x)=g(x)$ 的实数解个数，并说明它等于哪两个点集的交点个数；
(2) 讨论方程 $f(x)=a$ 的实数解个数随参数 $a$ 的变化情况。
\end{problem}
\begin{solution}
\paragraph{集合与点集本质剖析}
方程 $f(x)=g(x)$ 的解集是 $\{x\mid f(x)=g(x)\}$，几何上正是曲线 $y=f(x)$ 与 $y=g(x)$ 交点的横坐标集合，即两个点集的交集在 $x$ 轴上的投影。而 $f(x)=a$ 则是抛物线与水平线 $y=a$ 的交点问题，把“解的个数”转化为“点集交集元素个数”。
\paragraph{标准解题步骤}
(1) $x^2-2x=x\Rightarrow x^2-3x=0\Rightarrow x(x-3)=0$，得 $x=0$ 或 $x=3$，对应交点 $(0,0)$、$(3,3)$，故有 $2$ 个实数解。(2) 方程化为 $x^2-2x-a=0$，判别式 $\Delta=4+4a=4(1+a)$。当 $a<-1$ 时 $\Delta<0$，无实数解，几何上直线 $y=a$ 在抛物线顶点下方；当 $a=-1$ 时 $\Delta=0$，有 $1$ 个解 $x=1$，直线与抛物线相切；当 $a>-1$ 时 $\Delta>0$，有 $2$ 个解 $x=1\pm\sqrt{1+a}$。
\paragraph{学生典型易错点诊断}
最易犯的错误是不做判别式讨论，直接说“$\Delta>0$ 所以两解”，漏掉 $\Delta\leq0$ 的情形；二是把 $\Delta=0$ 误判为无解，其实此时是切点、恰有一解；三是几何与代数脱节，不画图导致对 $a=-1$ 的临界情形没有感觉；四是写出求根公式时忘记 $x=1\pm\sqrt{1+a}$ 只在 $a>-1$ 时有意义。规避方法：先写判别式、再按 $\Delta$ 的符号分类，最后回到图象核对交点个数。
\end{solution}

\begin{problem}\level{中}
设 $f(x)=x^2$，$g(x)=x+2$。
(1) 解不等式 $f(x)>g(x)$；
(2) 结合图象说明解集的几何意义，并指出边界点的位置。
\end{problem}
\begin{solution}
\paragraph{集合与点集本质剖析}
$f(x)>g(x)$ 的几何意义是：在同一竖直线 $x=a$ 上，点 $(a,f(a))$ 在点 $(a,g(a))$ 的上方。于是不等式的解集就是“抛物线位于直线上方”时 $x$ 的集合，本质仍是两个点集的上下位置关系。
\paragraph{标准解题步骤}
$x^2>x+2\Rightarrow x^2-x-2>0\Rightarrow (x-2)(x+1)>0$，解得 $x<-1$ 或 $x>2$，解集为 $(-\infty,-1)\cup(2,+\infty)$。两曲线交点由 $x^2=x+2$ 得 $x=-1,2$，交点为 $(-1,1)$、$(2,4)$，它们正是解集的边界点；在 $(-1,2)$ 内直线在抛物线上方，不等式不成立。对应的平面区域可表示为 $\{(x,y)\mid x<-1\text{ 或 }x>2,\ x+2<y<x^2\}$。
\paragraph{学生典型易错点诊断}
最易犯的错误是对 $x^2>x+2$ 两边直接除以 $x$，因未讨论 $x$ 的符号而得出错误解集；二是忘记边界点 $-1,2$ 处取等号，原不等式是严格大于，解集不含端点；三是把解集方向写反成 $[-1,2]$，或误写成闭区间；四是求出交点后不画图，无法判断哪一段抛物线在上方。规避方法：先求交点把数轴分段，再在每段判断上下位置，最后按严格不等号取舍端点。
\end{solution}

\begin{problem}\level{中}
画出不等式组 $\begin{cases}x\geq0,\\ y\geq0,\\ x+y\leq2,\\ x\leq y\end{cases}$ 表示的平面区域，判断各边界线应画实线还是虚线，并求该区域的面积。
\end{problem}
\begin{solution}
\paragraph{集合与点集本质剖析}
每个二元不等式确定半个平面，不等式组的解集是这些半平面的交集，即点集 $D=\{(x,y)\mid x\ge0,\ y\ge0,\ x+y\le2,\ x\le y\}$。交集通常是一个多边形区域（可行域），边界由直线段组成，这是线性规划中“可行域”的雏形。
\paragraph{标准解题步骤}
四条边界直线为 $x=0$、$y=0$、$x+y=2$、$y=x$。由于四个不等式都不带严格不等号，区域包含全部边界，四条边界线均画实线。求顶点：$x=0$ 与 $y=x$ 交于 $(0,0)$；$x=0$ 与 $x+y=2$ 交于 $(0,2)$；$y=x$ 与 $x+y=2$ 交于 $(1,1)$。区域是三角形，三个顶点为 $(0,0)$、$(0,2)$、$(1,1)$。以 $(0,0)$、$(0,2)$ 为底边，底长 $2$，点 $(1,1)$ 到 $y$ 轴的距离为高 $1$，故面积 $S=\dfrac12\times2\times1=1$。
\paragraph{学生典型易错点诊断}
最易犯的错误是把边界线画成虚线，误以为不等式组区域不含边界；二是忘记 $x\geq0,y\geq0$ 已把区域限制在第一象限，从而把区域误判为四边形；三是求交点时漏掉 $y=x$ 与 $x+y=2$ 的交点 $(1,1)$；四是算面积时用错底和高。规避方法：先逐条画边界线并判断虚实，再两两求交点确定顶点，最后用多边形面积公式（或分割法）计算。
\end{solution}

\section{深化·基本初等函数结构与逆关系}

\begin{problem}\level{难}
设 $f(x)=2^x\ (x\in\mathbb{R})$，其点集为 $G=\{(x,2^x)\mid x\in\mathbb{R}\}$。
(1) 写出反函数 $f^{-1}$ 的解析式与其点集；
(2) 说明 $G$ 与 $f^{-1}$ 的点集关于直线 $y=x$ 对称，并举例验证；
(3) 说明 $y=x^2$ 在 $\mathbb{R}$ 上没有反函数，而在 $[0,+\infty)$ 上却有。
\end{problem}
\begin{solution}
\paragraph{集合与点集本质剖析}
反函数的集合本质是把每个有序数对的两个坐标互换：$f^{-1}=\{(y,x)\mid (x,y)\in f\}$。几何上，这正好是关于直线 $y=x$ 的对称点集。指数函数与对数函数互为反函数，正体现了“指数结构”与“对数结构”的互逆。
\paragraph{标准解题步骤}
(1) 由 $y=2^x$ 解出 $x=\log_2 y$，故 $f^{-1}(x)=\log_2 x$，定义域 $(0,+\infty)$，值域 $\mathbb{R}$；其点集为 $G^{-1}=\{(2^x,x)\mid x\in\mathbb{R}\}=\{(x,\log_2 x)\mid x>0\}$。(2) 点 $(2,4)\in G$，关于 $y=x$ 的对称点 $(4,2)\in G^{-1}$，因为 $\log_2 4=2$。(3) $y=x^2$ 在 $\mathbb{R}$ 上，点 $(2,4)$ 与 $(-2,4)$ 都在图象上，互换后 $(4,2)$、$(4,-2)$ 使 $x=4$ 对应两个 $y$，破坏单值性；限制在 $[0,+\infty)$ 后每个正 $y$ 恰有一个原像，反函数为 $y=\sqrt{x}$。
\paragraph{学生典型易错点诊断}
最易犯的错误是把反函数理解成简单的“解出 $x$”，而忽略定义域与值域必须互换，把 $f^{-1}$ 的定义域写成 $\mathbb{R}$；二是认为 $y=x^2$ 总有反函数，忘记它不是一一映射；三是把对称轴写成 $y=-x$ 或 $x$ 轴；四是只验证一个点就下结论。规避方法：先写出互换后的点集，再检查“每个 $x$ 是否唯一”，最后核对定义域与值域互换。
\end{solution}

\begin{problem}\level{难}
设函数 $f$ 满足下列抽象方程之一，判断其原型（正比例函数、指数函数、对数函数），并求值：
(1) $f(x+y)=f(x)+f(y)$，$f(1)=2$，求 $f(3)$；
(2) $f(x+y)=f(x)f(y)$，$f(1)=2$，求 $f(3)$；
(3) $f(xy)=f(x)+f(y)\ (x,y>0)$，$f(2)=1$，求 $f(8)$。
\end{problem}
\begin{solution}
\paragraph{集合与点集本质剖析}
三个抽象方程分别对应三类基本初等函数的结构特征：加法保持不变的是正比例函数 $f(x)=kx$；加法变成乘法的是指数函数 $f(x)=a^x$；乘法变成加法的是对数函数 $f(x)=\log_a x$。识别原型后，求值就是套用该函数的结构性质。
\paragraph{标准解题步骤}
(1) 原型为 $f(x)=kx$，由 $f(1)=2$ 得 $k=2$，故 $f(3)=6$；也可直接写 $f(3)=f(1+1+1)=3f(1)=6$。(2) 原型为 $f(x)=a^x$，由 $f(1)=2$ 得 $a=2$，故 $f(3)=2^3=8$；也可写 $f(3)=f(1)^3=8$。(3) 原型为 $f(x)=\log_a x$，由 $f(2)=1$ 得 $a=2$，故 $f(8)=\log_2 8=3$；也可写 $f(8)=f(2\cdot2\cdot2)=3f(2)=3$。
\paragraph{学生典型易错点诊断}
最易犯的错误是把三个方程的原型互相混淆，例如把 $f(x+y)=f(x)f(y)$ 当成幂函数或正比例函数；二是忽视 (3) 中 $x,y>0$ 的定义域限制，把 $\log_a x$ 当成对任意实数有定义；三是直接假设 $f(x)=x$ 之类具体函数代入，却不说明依据；四是把 (2) 的结果写成 $f(3)=3f(1)$ 而误得 $6$。规避方法：先背熟“加法保持→正比例、加变乘→指数、乘变加→对数”，再用原型函数求值。
\end{solution}

\begin{problem}\level{难}
设 $f$ 满足 $f(x_1+x_2)=f(x_1)f(x_2)$，且 $f(1)=3$；设 $g$ 满足 $g(x_1x_2)=g(x_1)+g(x_2)\ (x_1,x_2>0)$，且 $g(3)=1$。
(1) 求 $f(4)$ 与 $g(9)$；
(2) 指出这两个方程分别“特化”了哪一条指数或对数运算律。
\end{problem}
\begin{solution}
\paragraph{集合与点集本质剖析}
抽象函数方程本质上是运算律的“结构特化”：$f(x_1+x_2)=f(x_1)f(x_2)$ 正是 $a^{x_1+x_2}=a^{x_1}a^{x_2}$，$g(x_1x_2)=g(x_1)+g(x_2)$ 正是 $\log_a(x_1x_2)=\log_a x_1+\log_a x_2$。结构相同，就可以借助原型函数完成求值。
\paragraph{标准解题步骤}
(1) 对 $f$：$f(4)=f(1+1+1+1)=f(1)^4=3^4=81$，也可由原型 $f(x)=3^x$ 得 $f(4)=3^4=81$。对 $g$：$g(9)=g(3\cdot3)=g(3)+g(3)=2$，也可由原型 $g(x)=\log_3 x$ 得 $g(9)=\log_3 9=2$。(2) 前者特化指数运算法则 $a^{m+n}=a^m a^n$，后者特化对数运算法则 $\log_a(MN)=\log_a M+\log_a N$。
\paragraph{学生典型易错点诊断}
最易犯的错误是把 $f(4)$ 算成 $4f(1)=12$，混淆了“加法型”与“乘法型”结构；二是把 $g(9)$ 算成 $g(3)^2=1$，同样是不理解“乘变加”的含义；三是忽略 $g$ 只在正实数上有定义，随意代入负数或零；四是只写答案不指出对应的运算律，导致“结构特化”这一核心观点没有落地。规避方法：先辨认方程属于哪一种运算结构，再调用对应的指数或对数法则。
\end{solution}

\section{拓展·反例与边界}

\begin{problem}\level{拓}
判断下列分段对应是不是函数；若是，求定义域与值域；若不是，指出破坏单值性的点。
(1) $f(x)=\begin{cases}x^2,&x<1,\\2x-1,&x>1;\end{cases}$
(2) $g(x)=\begin{cases}x+1,&x\leq1,\\2x,&x\geq1;\end{cases}$
(3) $h(x)=\begin{cases}x+1,&x\leq1,\\2x+1,&x\geq1.\end{cases}$
\end{problem}
\begin{solution}
\paragraph{集合与点集本质剖析}
分段函数的点集是各段点集的并集，因此要检查两件事：各段自身的定义域是否留下空缺，以及分段点处各段取值是否一致。只有“每个 $x$ 都恰好取到一个值”时，并集才满足单值性。
\paragraph{标准解题步骤}
(1) 两段在 $x=1$ 处都无定义，故定义域为 $(-\infty,1)\cup(1,+\infty)$；两段定义域不交，各自单值，所以 $f$ 是函数，值域为 $[0,+\infty)$（$x<1$ 时 $x^2$ 已取遍非负实数）。(2) $x=1$ 处两段分别给出 $2$ 与 $2$，取值一致，故 $g$ 是函数，定义域为 $\mathbb{R}$，$g(1)=2$。(3) $x=1$ 处 $x+1=2$ 而 $2x+1=3$，两值矛盾，破坏单值性，故 $h$ 不是函数。
\paragraph{学生典型易错点诊断}
最易犯的错误是看到“分段”就断言不是函数，其实分段函数多数是函数；二是忽略分段点处两段同时有定义的情形，不检查取值是否一致；三是求 (1) 的定义域时误写成 $\mathbb{R}$，漏掉 $x=1$ 这个空缺；四是求值域时漏掉分段点处的值。规避方法：先求各段定义域的并集，再单独检查分段点，最后逐段求值域。
\end{solution}

\begin{problem}\level{拓}
狄利克雷函数定义为 $D(x)=\begin{cases}1,&x\in\mathbb{Q},\\0,&x\notin\mathbb{Q}.\end{cases}$
(1) 求 $D(0)$、$D\left(\dfrac12\right)$、$D(\sqrt2)$；
(2) 判断 $D(x)$ 的奇偶性；
(3) 判断“对任意非零有理数 $T$，都有 $D(x+T)=D(x)$”，并说明理由；
(4) 说明为什么 $D(x)$ 的图象“画不出来”。
\end{problem}
\begin{solution}
\paragraph{集合与点集本质剖析}
$D$ 的点集被拆成两个稠密但交错的子集：$G_1=\{(x,1)\mid x\in\mathbb{Q}\}$ 与 $G_0=\{(x,0)\mid x\notin\mathbb{Q}\}$。有理点集与无理点集在 $\mathbb{R}$ 上都稠密，任何区间内两类点都会出现，这正是全部“怪异”性质的来源。
\paragraph{标准解题步骤}
(1) $0\in\mathbb{Q}$，$\dfrac12\in\mathbb{Q}$，故 $D(0)=D\left(\dfrac12\right)=1$；$\sqrt2\notin\mathbb{Q}$，故 $D(\sqrt2)=0$。(2) 对任意 $x$，$x$ 与 $-x$ 同为有理或同为无理，故 $D(-x)=D(x)$，$D$ 是偶函数。(3) 若 $T\in\mathbb{Q}\setminus\{0\}$，则 $x+T$ 与 $x$ 同为有理或同为无理，故 $D(x+T)=D(x)$；任何非零有理数都是周期，且不存在最小正周期。(4) 任何区间内图象都无限密集地分布在 $y=0$ 与 $y=1$ 两条直线上，无法用有限笔画出。
\paragraph{学生典型易错点诊断}
最易犯的错误是认为 $D(x)$ 不是周期函数，其实每个非零有理数都是它的周期；二是把 $\sqrt2$ 当成有理数；三是把“画不出图象”误解为“函数不存在”，混淆了“图象不可视”与“对应不存在”；四是判断奇偶性时忽略有理数集的对称性。规避方法：时刻抓住“有理数与无理数”这一分类，先判断自变量属于哪一类，再代入求值。
\end{solution}

\begin{problem}\level{拓}
已知圆 $x^2+y^2=4$。
(1) 用垂线检验法说明圆不是函数图象；
(2) 把下半圆写成函数解析式，并求其定义域与值域；
(3) 说明为什么在 $x=\pm2$ 处仍保持单值，并求下半圆在 $x=0$ 处的函数值。
\end{problem}
\begin{solution}
\paragraph{集合与点集本质剖析}
圆上的点集是 $\{(x,y)\mid x^2+y^2=4\}$。对 $x\in(-2,2)$，同一竖直线与圆交于上下两点，破坏单值性；把圆按 $x$ 轴拆成上半圆与下半圆，每一段都满足垂线检验，这就是“局部单值化”。
\paragraph{标准解题步骤}
(1) 取 $x=0$，圆上有 $(0,2)$ 与 $(0,-2)$ 两点，竖直线 $x=0$ 与圆有两个交点，故整个圆不是函数图象。(2) 下半圆为 $y=-\sqrt{4-x^2}$，由偶次根号内非负得 $4-x^2\geq0$，即 $-2\leq x\leq2$，定义域为 $[-2,2]$；因为 $0\leq\sqrt{4-x^2}\leq2$，所以值域为 $[-2,0]$。(3) 在 $x=\pm2$ 处只有唯一一点 $(\pm2,0)$，竖直线只与图象交于一点，故端点处仍单值；$x=0$ 时 $y=-\sqrt{4}=-2$，即下半圆的最低点 $(0,-2)$。
\paragraph{学生典型易错点诊断}
最易犯的错误是把下半圆的定义域写成开区间 $(-2,2)$，丢掉两个端点；二是忘记 $\sqrt{4-x^2}$ 非负，误以为 $y=-\sqrt{4-x^2}$ 的值域是 $[0,2]$，符号判断出错；三是把上半圆与下半圆混为一谈；四是认为“圆不是函数”就可以随意写成 $y=\pm\sqrt{4-x^2}$ 并称之为函数。规避方法：先写解析式并求定义域，再根据根号的非负性确定值域，最后用垂线检验逐段核对。
\end{solution}
